\scanpage{32}

$\Bigg[$ 1. Généralisation des th. de finitude formelle, et des th. de comparaison. Application : un th. d'existence

2. Le théorème de finitude locale et ses applications.

\deleted{3. Faisceaux de Lefschetz-Grauert. Remarques.}

\deleted{4. \unsure{Propriétés} propriétés \unsure{des} \uline{Exemples} \unsure{fondament}} \inserted{1}

3. Cas particulier 1 : \deleted{\unsure{Etude}} \inserted{théorie} des \deleted{\unsure{spa}} schémas locaux.

4. Cas particulier 2 : \deleted{\unsure{Etude}} \inserted{théorie} des schémas \inserted{algébriques} projectifs.

5. Faisceaux de Lefschetz-Grauert : Remarques, exemples.

% slanted
\margin{utilise IV, \S\ sur les morphismes étales....}
6. Propriétés liées aux $\pi_0$, \tikz[baseline=(p.base)]\node[draw,ellipse,inner sep=1pt](p){$\pi_1$};, Pic.

[7. Propriétés de permanence]

\textcircled{8}. Théorèmes de pureté, et théorèmes de Lefschetz pour le $\pi_1$.

\deleted{\textcircled{\unsure{9|1}}} 9 \uline{Anneaux} \deleted{\unsure{locaux}} \uline{quasi-factoriels} et \uline{th. de Lefschetz pour Pic}.

% archive sheet number: 31
